% FILE: Supplement protocol guide — supplement list, dosages, protocols, recommendations
\chapter{Practical Supplement Guide}
\label{app:supplement-guide}

This appendix provides a practical reference for supplements commonly used in ME/CFS management, organized by target mechanism. All information is derived from the evidence reviewed in Chapter~\ref{ch:supplements}. This guide is intended as a reference for informed discussion between patients and their healthcare providers, not as a prescription.

\begin{warning}[Medical Supervision Required]
Supplements can interact with medications, affect laboratory results, and produce adverse effects. All supplement decisions should be made in consultation with a qualified healthcare provider who is aware of the patient's complete medical history and medication list. The dosages listed below are derived from published studies and clinical protocols; individual requirements may differ.
\end{warning}

\section{Quick Reference Tables}
\label{sec:supplement-tables}

\subsection{Mitochondrial and Energy Support}

Table~\ref{tab:mito-supplements} lists supplements targeting energy metabolism dysfunction, the primary pathophysiological mechanism in ME/CFS (Chapter~\ref{ch:energy-metabolism}).

\begin{table}[htbp]
\centering
\caption{Mitochondrial and energy support supplements.}
\label{tab:mito-supplements}
\small
\begin{tabular}{p{2.8cm}p{2.5cm}p{1.5cm}p{4.5cm}}
\hline
\textbf{Supplement} & \textbf{Typical Dose} & \textbf{Evidence} & \textbf{Mechanism and Notes} \\
\hline
Coenzyme Q10 (ubiquinol) & 200--400\,mg/day & Moderate & ETC electron carrier; RCT showed reduced fatigue~\cite{CastroMarrero2021CoQ10}. Take with fat-containing meal. \\
NAD$^+$ precursors (NR or NMN) & 250--500\,mg/day & Low & Replenishes NAD$^+$ pool; limited ME/CFS-specific data~\cite{syed2025nad_therapy}. \\
D-Ribose & 5\,g 2--3$\times$/day & Low & ATP precursor; small open-label studies suggest benefit. Dissolve in water. \\
Acetyl-L-carnitine & 1000--2000\,mg/day & Low & Fatty acid transport into mitochondria; supports $\beta$-oxidation. \\
Alpha-lipoic acid & 300--600\,mg/day & Low & Antioxidant; Krebs cycle cofactor. R-lipoic acid form preferred. \\
B-complex vitamins & 1$\times$/day (balanced formula) & Low & Cofactors for energy metabolism; B2 (riboflavin) for Complex~I, B3 (niacin) for NAD$^+$. \\
\hline
\end{tabular}
\end{table}

\subsection{Immune Modulation}

\begin{table}[htbp]
\centering
\caption{Immune modulation supplements.}
\label{tab:immune-supplements}
\small
\begin{tabular}{p{2.8cm}p{2.5cm}p{1.5cm}p{4.5cm}}
\hline
\textbf{Supplement} & \textbf{Typical Dose} & \textbf{Evidence} & \textbf{Mechanism and Notes} \\
\hline
Vitamin D3 & 2000--5000\,IU/day & Moderate & Immunomodulatory; deficiency common in housebound patients. Monitor serum levels; target 40--60\,ng/mL. \\
Vitamin C & 500--1000\,mg/day & Low & Antioxidant; immune cell support. Higher doses may cause GI upset. \\
Zinc & 15--30\,mg/day & Low & NK cell function; thymic hormone production. Long-term use requires copper co-supplementation. \\
Selenium & 100--200\,$\mu$g/day & Low & Antioxidant selenoproteins; thyroid function. Do not exceed 400\,$\mu$g/day. \\
Quercetin & 500--1000\,mg/day & Low & Mast cell stabilizer; anti-inflammatory. May help MCAS-overlap symptoms. \\
\hline
\end{tabular}
\end{table}

\subsection{Neurological and Sleep Support}

\begin{table}[htbp]
\centering
\caption{Neurological and sleep support supplements.}
\label{tab:neuro-supplements}
\small
\begin{tabular}{p{2.8cm}p{2.5cm}p{1.5cm}p{4.5cm}}
\hline
\textbf{Supplement} & \textbf{Typical Dose} & \textbf{Evidence} & \textbf{Mechanism and Notes} \\
\hline
Magnesium (glycinate or threonate) & 200--400\,mg elemental/day & Low & Muscle relaxation; NMDA receptor modulation; sleep support. Glycinate for general use; threonate for cognitive symptoms. \\
Melatonin & 0.5--3\,mg at bedtime & Low & Circadian rhythm regulation~\cite{castromarrero2021melatonin}. Start low; higher doses not more effective. \\
Omega-3 fatty acids (EPA/DHA) & 1000--2000\,mg combined/day & Low & Anti-inflammatory; neuronal membrane support. Choose high-EPA formulations for mood. \\
L-Theanine & 100--200\,mg/day & Very low & GABAergic; promotes relaxation without sedation. \\
\hline
\end{tabular}
\end{table}

\subsection{Gut Health}

\begin{table}[htbp]
\centering
\caption{Gut health supplements.}
\label{tab:gut-supplements}
\small
\begin{tabular}{p{2.8cm}p{2.5cm}p{1.5cm}p{4.5cm}}
\hline
\textbf{Supplement} & \textbf{Typical Dose} & \textbf{Evidence} & \textbf{Mechanism and Notes} \\
\hline
Probiotics (multi-strain) & $10^9$--$10^{10}$ CFU/day & Low & Microbiome modulation; strain-specific effects~\cite{Hsu2025gut}. Lactobacillus and Bifidobacterium strains most studied. \\
Butyrate (tributyrin or sodium butyrate) & 300--600\,mg 2--3$\times$/day & Very low & Intestinal barrier support; anti-inflammatory~\cite{Folkerts2020butyrate}. Enteric-coated preferred. \\
L-Glutamine & 5--10\,g/day & Very low & Intestinal epithelial cell fuel; barrier integrity. \\
\hline
\end{tabular}
\end{table}

Evidence levels are defined as: \textbf{Moderate} = at least one RCT in ME/CFS population; \textbf{Low} = open-label studies, case series, or extrapolation from related conditions; \textbf{Very low} = mechanistic rationale only, no ME/CFS-specific clinical data.

\section{Supplement Interactions}
\label{sec:supplement-interactions}

\subsection{Supplement--Supplement Interactions}

Most supplements listed above can be taken concurrently without adverse interactions. Notable considerations:

\begin{itemize}
    \item \textbf{Zinc and copper}: Long-term zinc supplementation ($>30$\,mg/day for $>8$ weeks) can induce copper deficiency. Co-supplement with 1--2\,mg copper if zinc exceeds 30\,mg/day.
    \item \textbf{CoQ10 and blood thinners}: CoQ10 has structural similarity to vitamin~K and may reduce warfarin efficacy. Monitor INR if taking both.
    \item \textbf{Iron and other minerals}: Iron inhibits absorption of zinc, calcium, and magnesium. Separate iron from other mineral supplements by $\geq 2$ hours.
    \item \textbf{Alpha-lipoic acid and thyroid medication}: May reduce levothyroxine absorption. Separate by $\geq 4$ hours.
    \item \textbf{Melatonin and immunosuppressants}: Melatonin has immunomodulatory properties; use with caution in patients on immunosuppressive therapy.
\end{itemize}

\subsection{Supplement--Drug Interactions}

Interactions with commonly prescribed ME/CFS medications:

\begin{itemize}
    \item \textbf{Low-dose naltrexone (LDN)}: No known interactions with the supplements listed above.
    \item \textbf{Beta-blockers}: CoQ10 may modestly reduce blood pressure; monitor if combining with propranolol or other antihypertensives.
    \item \textbf{SSRIs/SNRIs}: Omega-3 fatty acids may potentiate anticoagulant effects; 5-HTP (not listed above) should be avoided with serotonergic medications due to serotonin syndrome risk.
    \item \textbf{Fludrocortisone}: Magnesium supplementation is advisable as fludrocortisone can increase magnesium excretion.
    \item \textbf{Midodrine}: No significant interactions with listed supplements.
\end{itemize}

\section{Quality and Sourcing}
\label{sec:quality-sourcing}

Supplement quality varies substantially because dietary supplements are not subject to the same regulatory standards as pharmaceuticals. Patients should prioritize products verified by independent testing organizations:

\begin{itemize}
    \item \textbf{USP (United States Pharmacopeia)}: Tests for identity, potency, purity, and dissolution
    \item \textbf{NSF International}: Certifies Good Manufacturing Practices compliance and contaminant testing
    \item \textbf{ConsumerLab}: Independent testing with publicly available results
    \item \textbf{Informed Sport/Informed Choice}: Tests for banned substances (relevant for athletes)
\end{itemize}

Red flags suggesting poor quality:

\begin{itemize}
    \item Claims of ``curing'' or ``treating'' ME/CFS
    \item Proprietary blends without individual ingredient quantities
    \item Lack of third-party testing certification
    \item Unrealistically low prices for expensive ingredients (e.g., ubiquinol, NMN)
    \item Marketing that relies on testimonials rather than cited research
\end{itemize}

\section{Starting a Supplement Protocol}
\label{sec:starting-protocol}

ME/CFS patients are frequently sensitive to new substances, including supplements. A systematic introduction protocol reduces the risk of adverse reactions and enables identification of which supplements provide benefit:

\begin{enumerate}
    \item \textbf{One supplement at a time}: Introduce each new supplement individually, with a minimum 1--2 week observation period before adding the next. This ensures that any reaction (positive or negative) can be attributed to the correct supplement.
    \item \textbf{Start low}: Begin at one-quarter to one-half of the target dose and increase gradually over 1--2 weeks if tolerated. ME/CFS patients may respond to lower doses than the general population.
    \item \textbf{Priority order}: Start with supplements that have the strongest evidence and target the patient's dominant symptoms:
    \begin{itemize}
        \item Fatigue-dominant: CoQ10, then NAD$^+$ precursors, then D-ribose
        \item Sleep-dominant: Magnesium glycinate, then melatonin
        \item Immune/pain-dominant: Vitamin D (if deficient), then omega-3, then quercetin
        \item GI-dominant: Probiotics, then butyrate
    \end{itemize}
    \item \textbf{Track responses}: Maintain a simple log of symptoms before and after each supplement introduction. A minimum 4-week trial at target dose is needed to assess efficacy; some supplements (especially CoQ10, vitamin D) require 8--12 weeks.
    \item \textbf{Periodic reassessment}: Every 3--6 months, consider whether each supplement is still providing benefit. Discontinue one supplement at a time and observe for 2--4 weeks. If no worsening occurs, the supplement may be unnecessary.
\end{enumerate}

\section{When to Stop a Supplement}
\label{sec:when-to-stop}

Discontinue a supplement and consult a healthcare provider if:

\begin{itemize}
    \item New symptoms develop that correlate temporally with supplement introduction
    \item Gastrointestinal symptoms (nausea, diarrhea, cramping) persist beyond the first week
    \item Allergic reaction signs appear (rash, itching, swelling, breathing difficulty)
    \item No discernible benefit after an adequate trial period (4--12 weeks depending on supplement)
    \item Laboratory values move outside reference ranges (e.g., vitamin~D $>80$\,ng/mL, selenium $>200$\,$\mu$g/L)
    \item Cost becomes prohibitive relative to perceived benefit
\end{itemize}

\begin{limitation}[Evidence Gaps in ME/CFS Supplementation]
Most supplements used in ME/CFS management lack randomized controlled trial evidence specific to ME/CFS populations. The recommendations in this appendix are based on a combination of limited ME/CFS-specific data, mechanistic rationale from the pathophysiology described in Part~\ref{part:pathophysiology}, and extrapolation from related conditions (fibromyalgia, mitochondrial diseases, post-viral fatigue). Patients should be informed that supplementation is empirical and that individual responses are unpredictable.
\end{limitation}
